% chap1.tex
\part{서론}
\label{part:introduction}

% ---------------------------------------------------------------------------- %
%                                  NEW SECTION                                 %
% ---------------------------------------------------------------------------- % 

\chapter{서론}

\section{개발 배경 및 목적}

우리 정부는 2030 국가온실가스감축목표(NDC2030)를 통해 중기 감축 경로를 제시하고, 2050 탄소중립을 국가 비전으로 채택함으로써 건물 부문의 에너지 소비와 온실가스 배출을 구조적으로 줄여 나가야 하는 과제를 안고 있다\cite{korea_ndc_2050cnc_website}\cite{korea_ndc_enhanced_2021}.
이러한 목표 달성 과정에서 건물 부문은 제로에너지건축물(ZEB) 의무화 정책을 통해 신축 건물에서 고효율·저에너지 구조를 선제적으로 확보하는 한편, 이미 존재하는 방대한 기축 건축물에 대해서는 외피 성능 개선, 고효율 설비 도입, 재생에너지 활용 등을 종합적으로 검토하는 체계적인 그린리모델링을 통해 성능을 단계적으로 향상시키는 전략을 취하고 있다.

이 때, 개별 건물의 그린리모델링은 단열 강화, 창호 교체, 고효율 냉난방 설비 도입, 환기 및 제어 시스템 개선, 재생에너지 설비 연계 등에 걸쳐 복수의 요소기술 조합을 채택할 수 있으며, 각 대안은 초기 투자비, 유지관리 비용, 공사 난이도, 입주자 불편, 법·제도 요건에 이르기까지 서로 다른 특성을 가진다. 또한, 동일한 요소기술 조합이라도 건물의 용도, 규모, 기후 조건에 따라 효과와 경제성이 크게 달라지기 때문에\cite{shen2019feasibility}, 그린리모델링 사업의 의사결정은 일률적으로 이루어질 수 없고 개별건물의 특성에 맞게 적용되어야 의도한 에너지 절감 효과를 극대화할 수 있다.
결국 그린리모델링 의사결정의 핵심은 개별 건물의 맥락과 제약 조건 속에서 여러 요소기술 조합을 비교하여 해당 건물에 가장 적합한 대안을 선택하는 것이며, 이를 위해서는 각 대안이 시간에 따라 건물 거동과 에너지 흐름에 어떤 영향을 미치는지를 사전에 시뮬레이션 할 수 있는 도구가 필수적이다\cite{chen2017automatic}\cite{mostafazadeh2023energy}.

이러한 맥락에서 \textbf{동적 건물 에너지 시뮬레이션(dynamic building energy simulation)}은 시간에 따라 변하는 건물의 상태와 에너지 흐름을 정교하게 설명할 수 있는 방법으로서, 그린리모델링을 위한 분석,설계, 및 평가 전 과정에서 높은 잠재적 가치를 가진다\cite{kotarela2023exploitation}\cite{lops2025towards}.
동적 시뮬레이션은 기상 변화(외기온, 일사, 풍속 등), 내부부하(인체, 조명, 기기, 침기, 환기), 구조체의 열적 거동(축열 및 방열), 및 설비시스템의 거동(ON/OFF 제어, 부분부하 등)을 짧은 시간 간격(1시간 이내, 보통 10분)으로 계산하여, 건물의 내부 온도, 열부하, 설비 소비전력, 실내 열쾌적 상태를 구하며, 이를 통해 연간 에너지 소비량, 피크 부하, 온실가스 배출량, 실내 쾌적성 지표 등을 높은 정밀도로 상세히 산출할 수 있다.
특히 그린리모델링 이후의, 고단열,고축열 저-에너지 건축물에서는, 이러한 건물의 동특성이 연간 부하에 미치는 영향이 더욱 강조된다 \parencref{fig:vandijk-whydynamic}.

\begin{defaultfigure}
  \centering
  
  \begin{subfigure}[b]{0.48\textwidth}
    \centering
    \includegraphics[width=\textwidth]{VanDijk2018다이어그램-과거건물.png}
    \subcaption{과거의 저단열, 저축열 건축물}
  \end{subfigure}
  \hfill
  \begin{subfigure}[b]{0.48\textwidth}
    \centering
    \includegraphics[width=\textwidth]{VanDijk2018다이어그램-저에너지건물.png}
    \subcaption{현대의 고단열, 고축열 건축물}
  \end{subfigure}
  
  \caption{과거 건축물과 그린리모델링 이후 저에너지 건축물의 동적 거동 차이\cite{van2018epb}}
  \label{fig:vandijk-whydynamic}
\end{defaultfigure}

그럼에도 불구하고, 이러한 동적 시뮬레이션의 장점은 아래와 같은 이유로 국내 그린리모델링 실무에서 충분히 활용되고 있다고 보기는 어렵다.
\begin{itemize}
  \item 첫째, 동적 시뮬레이션은 \textbf{입력 데이터의 구조와 양이 상대적으로 방대}하다. 존 구획, 외피 구성, 창호 성능, 내부 발열, 설비 시스템 구성, 제어 로직, 운전 스케줄 등 다양한 요소를 동시에 정의해야 하며, 각 요소가 시간에 따라 어떻게 변하는지까지 고려해야 한다. 이는 그린리모델링 사업에서 제한된 예산과 기간 안에 빠르게 분석을 수행해야 하는 실무자에게는 부담으로 작용할 수 있다.
  \item 둘째, 기존의 동적 시뮬레이션 도구들은 대체로 해외에서 개발된 범용 엔진이 많아, \textbf{인터페이스와 문서, 예제가 외국 기준과 용어 체계}에 맞추어져 있다. 이로 인해 국내 사용자들이 우리나라의 건축법규, 에너지 절약설계기준, 실제 설비 관행을 반영하여 활용하기 위해서는 추가적인 해석과 변환 작업이 필요하며, 초기 진입장벽이 높게 느껴진다.
  \item 셋째, 그린리모델링 사업에서는 연간 에너지 절감률, 투자 회수기간, 온실가스 감축량, 공사 전·후 쾌적성 개선 정도 등 정책·사업 관점의 지표가 중요하나, 기존 동적 시뮬레이션 도구들의 출력은 시간별 온도, 부하, 소비전력과 같은 시계열 데이터와 다차원 결과가 중심이기 때문에, 이를 \textbf{사업자가 이해하기 쉬운 지표와 그림으로 해석하는 절차가 요구}된다. 다시 말해, 동적 시뮬레이션 자체의 설명력은 매우 높지만, 그 설명력을 사업·정책 의사결정의 언어로 연결해 주는 인터페이스는 부족하다.
  \item 넷째, 국내에서는 동적 시뮬레이션을 활용한 그린리모델링 사례와 교육 자료, \textbf{표준화된 절차가 제한적}이다. 이는 도구의 신뢰성이 부족해서가 아니라, 축적된 사례와 레퍼런스가 충분히 공유되지 않았기 때문으로, 실제 프로젝트에서 어떤 수준의 모델링 정밀도가 적정한지, 어떤 입력은 단순화해도 괜찮고 어떤 부분은 반드시 동적으로 표현해야 하는지, 결과를 어떻게 검증하고 문서화할 것인지 등이 모두 실무자 입장에서 의사결정 리스크로 느껴진다.
\end{itemize}

이와 같은 요인들로 인해, 동적 시뮬레이션은 그린리모델링에 매우 적합한 도구임에도 불구하고, 국내에서는 소수의 전문가 집단에 의해 제한적으로 활용되는 경향이 있다. 그러나 건물의 저에너지화, 재생에너지 보급 확대, 동적 부하와 연계된 설비 운전 등, 그린리모델링의 핵심 이슈들이 본질적으로 시간적 변동성과 건축물의 다양한 요소 간 상호작용을 다루는 문제인 점을 감안하면, 그린리모델링 사업에서 동적 시뮬레이션의 필요성은 점점 더 증대됨이 명확하며, 국내 그린리모델링 환경에 특화된 동적 시뮬레이션 엔진이 필요하다.

\simulator 은, 이러한 문제와 필요성을 바탕으로, 아래의 목표를 달성하여 국내 그린리모델링 실무 환경에 적합한 동적 건물 에너지 시뮬레이션 엔진으로 기능하기 위해 개발되었다.

\begin{itemize}
  \item 첫째, \ep 를 엔진으로 활용하여, \textbf{국제적으로 검증된 high-fidelity 물리식 기반 시뮬레이션 알고리즘}으로 그린리모델링 전/후에 대한 동적 시뮬레이션을 수행할 수 있도록 한다.
  \item 둘째, 국내 그린리모델링 \textbf{실무 환경에 맞춘 입력 데이터 구조와 템플릿}을 제공하여, 사용자가 일반적인 동적 시뮬레이션 도구에 요구되는 복잡한  구조를 이해하지 않고도 분석이 가능하도록 한다.
  \item 셋째, \textbf{국내 건축법규, 에너지 절약설계기준, 그린리모델링 설계 가이드를 반영}한 기본값과 템플릿을 제공하여, 사용자가 별도의 변환 작업 없이도 국내 환경에 맞는 시뮬레이션을 수행할 수 있도록 한다.
  \item 넷째, 복잡한 동적 시뮬레이션 결과를 온실가스 배출량, 에너지비용 등 \textbf{그린리모델링 사업에 요구되는 지표}로 변환하여, 의사결정에 바로 활용할 수 있도록 한다.
  \item 다섯째, \textbf{시뮬레이션 알고리즘을 오픈소스로 공개}하여, 국내 연구자 및 실무자의 자유로운 검증으로 표준화된 절차에 대한 신뢰성을 높이고, 필요에 따라 기능을 확장 및 개량할 수 있도록 한다.
\end{itemize}

% ---------------------------------------------------------------------------- %
%                                  NEW SECTION                                 %
% ---------------------------------------------------------------------------- %

\section{\simulator 의 개념 및 구조}

\simulator 은 국제적으로 검증된 건물 에너지 해석 프로그램인 EnergyPlus(EP)를 핵심 연산 코어로 사용하되, 사용자가 EP의 복잡한 입출력 과정을 직접 다루지 않도록 설계된 추상화 계층(Abstraction Layer)이다 \parencref{fig:simulator_concept_structure}. 사용자는 그린리모델링 현장에서 수집 가능한 수준의 정보로 이루어진, tree 구조(json 형식)로 표준화된 건물 데이터 파일\footnote{이에 대한 상세한 기술은 \ref{section:grm-explanation} 참고}을 직접, 또는 적절한 GUI를 이용하여 생성하고, 이를 \simulator 에 전달한다. 이는 사용자가 복잡한 IDF 문법이나 수많은 상세 필드를 직접 다루지 않고도, 자신의 언어에 가까운 구조로 모델을 정의할 수 있음을 의미한다.

\begin{defaultfigure}
  \includegraphics[width=\textwidth]{시뮬레이터 개념 및 구조.png}
  \caption{\simulator 의 개념 및 구조}
  \label{fig:simulator_concept_structure}
\end{defaultfigure}

이후 Python 모듈 \pymodule\footnote{이에 대한 상세한 기술은 \ref{chapter:pymodule} 참고}이 이를 Python 객체 (그림에서, Input-Friendly class)로 변환하고, 데이터베이스 값 삽입, 세부 변수 가정, EnergyPlus 구조에 맞춘 재구성을 수행하여 EnergyPlus의 데이터 구조와 유사한 Python 객체(그림에서, EP-Friendly class)로 변환한다. 이는 템플릿 기반의 IDF 객체 생성 절차를 거쳐 idf 파일로 자동 변환되고, \ep 해석이 수행된 뒤 결과 파일이 다시 Python 객체 (그림에서, EnergyPlusResult) 객체로 읽혀 들어온다. 이후 이 결과는 월별·용도별 에너지 지표 등 그린리모델링 의사결정에 바로 활용 가능한 형태의 json 형식의 파일\footnote{이에 대한 상세한 기술은  \ref{section:grr-explanation} 참고}로 재가공된다.

즉, \simulator 은 입력 단계에서 건물 데이터 구조를 단순화·표준화하고, 해석 단계에서는 EnergyPlus의 물리 기반 고정밀 해석을 그대로 활용하며, 출력 단계에서는 실무자가 이해하기 쉬운 지표와 형식으로 재구성하는 Python 기반 변환 모듈이다. 이를 통해 사용자는 EnergyPlus의 복잡한 내부 구조를 직접 다루지 않고도, GUI와 단일 입출력 포맷만으로 고정밀 동적 시뮬레이션을 활용할 수 있게 된다.
